\documentclass[10pt,a4paper]{amsart}
\usepackage{amsmath,amssymb,amsthm,graphicx}
\usepackage{verbatim}
\newenvironment{reference}{\comment}{\endcomment}
%\newenvironment{reference}{}{}
\newenvironment{slogan}{\comment}{\endcomment}
\newenvironment{history}{\comment}{\endcomment}

% For commutative diagrams we use Xy-pic
\usepackage[all]{xy}

% We use 2cell for 2-commutative diagrams.
\xyoption{2cell}
% Theorem environments.
%
\theoremstyle{plain}
\newtheorem{theorem}[subsection]{Theorem}
\newtheorem{proposition}[subsection]{Proposition}
\newtheorem{lemma}[subsection]{Lemma}

\theoremstyle{definition}
\newtheorem{definition}[subsection]{Definition}
\newtheorem{example}[subsection]{Example}
\newtheorem{exercise}[subsection]{Exercise}
\newtheorem{situation}[subsection]{Situation}

\theoremstyle{remark}
\newtheorem{remark}[subsection]{Remark}
\newtheorem{remarks}[subsection]{Remarks}

\numberwithin{equation}{subsection}

% Macros
%
\def\lim{\mathop{\mathrm{lim}}\nolimits}
\def\colim{\mathop{\mathrm{colim}}\nolimits}
\def\Spec{\mathop{\mathrm{Spec}}}
\def\Hom{\mathop{\mathrm{Hom}}\nolimits}
\def\Ext{\mathop{\mathrm{Ext}}\nolimits}
\def\SheafHom{\mathop{\mathcal{H}\!\mathit{om}}\nolimits}
\def\SheafExt{\mathop{\mathcal{E}\!\mathit{xt}}\nolimits}
\def\Sch{\mathit{Sch}}
\def\Mor{\mathop{\mathrm{Mor}}\nolimits}
\def\Ob{\mathop{\mathrm{Ob}}\nolimits}
\def\Sh{\mathop{\mathit{Sh}}\nolimits}
\def\NL{\mathop{N\!L}\nolimits}
\def\CH{\mathop{\mathrm{CH}}\nolimits}
\def\proetale{{pro\text{-}\acute{e}tale}}
\def\etale{{\acute{e}tale}}
\def\QCoh{\mathit{QCoh}}
\def\Ker{\mathop{\mathrm{Ker}}}
\def\Im{\mathop{\mathrm{Im}}}
\def\Coker{\mathop{\mathrm{Coker}}}
\def\Coim{\mathop{\mathrm{Coim}}}

% Boxtimes
%
\DeclareMathSymbol{\boxtimes}{\mathbin}{AMSa}{"02}

%
% Macros for moduli stacks/spaces
%
\def\QCohstack{\mathcal{QC}\!\mathit{oh}}
\def\Cohstack{\mathcal{C}\!\mathit{oh}}
\def\Spacesstack{\mathcal{S}\!\mathit{paces}}
\def\Quotfunctor{\mathrm{Quot}}
\def\Hilbfunctor{\mathrm{Hilb}}
\def\Curvesstack{\mathcal{C}\!\mathit{urves}}
\def\Polarizedstack{\mathcal{P}\!\mathit{olarized}}
\def\Complexesstack{\mathcal{C}\!\mathit{omplexes}}
% \Pic is the operator that assigns to X its picard group, usage \Pic(X)
% \Picardstack_{X/B} denotes the Picard stack of X over B
% \Picardfunctor_{X/B} denotes the Picard functor of X over B
\def\Pic{\mathop{\mathrm{Pic}}\nolimits}
\def\Picardstack{\mathcal{P}\!\mathit{ic}}
\def\Picardfunctor{\mathrm{Pic}}
\def\Deformationcategory{\mathcal{D}\!\mathit{ef}}
\usepackage{cleveref}
\crefname{originalreference}{original reference}{original references}
\Crefname{originalreference}{Original reference}{Original references}
\makeatletter
\newlabel{lemma-base-change-tensor}{{0A1D}{source}{}{sourcepointer-1}{}}
\newlabel{lemma-base-change-tensor@cref}{{[originalreference][0][]0A1D}{[1][0][]source}{}{}{}}
\makeatother
\crefname{originalreference}{original reference}{original references}
\Crefname{originalreference}{Original reference}{Original references}
\makeatletter
\newlabel{lemma-base-change-tensor-perfect}{{0A1H}{source}{}{sourcepointer-2}{}}
\newlabel{lemma-base-change-tensor-perfect@cref}{{[originalreference][0][]0A1H}{[1][0][]source}{}{}{}}
\makeatother
\crefname{originalreference}{original reference}{original references}
\Crefname{originalreference}{Original reference}{Original references}
\makeatletter
\newlabel{lemma-quasi-coherence-direct-image}{{08D5}{source}{}{sourcepointer-3}{}}
\newlabel{lemma-quasi-coherence-direct-image@cref}{{[originalreference][0][]08D5}{[1][0][]source}{}{}{}}
\makeatother
\crefname{originalreference}{original reference}{original references}
\Crefname{originalreference}{Original reference}{Original references}
\makeatletter
\newlabel{theorem-approximation}{{08ES}{source}{}{sourcepointer-4}{}}
\newlabel{theorem-approximation@cref}{{[originalreference][0][]08ES}{[1][0][]source}{}{}{}}
\makeatother
\title{Pushforward of a pseudo-coherent object tensored with a nice complex is pseudo-coherent with base-change compatible formation}
\author{The Stacks Project Authors}
\date{stacks-project:spaces-perfect.tex:local-snapshot-20261012; GNU FDL 1.2 or later}
\begin{document}
\maketitle
\noindent\textit{Editorial note.} All delivered mathematical text is the author's own, copied from the pinned Stacks Project chapter file spaces-perfect.tex; no statement or proof step was rewritten, reordered or removed.\par\medskip
\noindent\textit{Editorial note.} The counted proof is the author's complete argument: the proof is complete: base change is quoted from the preceding lemma and m-pseudo-coherence is proved by approximation and truncation\par\medskip
\noindent\textit{Editorial note.} The label lemma-base-change-tensor-pseudo-coherent names the counted statement, and the proof environment immediately following it in the source is delivered verbatim as its proof.\par\medskip
\noindent\textit{Editorial note.} The counted statement is carried by the author's own lemma environment in the chapter source.\par\medskip
\noindent\textit{Allowed standard prerequisites (editorial).}\begin{itemize}
\item pseudo-coherent objects and m-pseudo-coherence
\item approximation of pseudo-coherent complexes by perfect complexes at a given threshold
\item cohomological dimension bound N for total direct image over quasi-compact quasi-separated spaces
\item the perfect base-change tensor lemma and the base-change compatibility of the total direct image of a tensor product
\item truncations and their interaction with derived tensor
\end{itemize}\medskip
\ifdefined\hypertarget\hypertarget{sourcepointer-1}{}\fi
\noindent\textit{Source reference.} Original reference 0A1D: sources/stacks-spaces-perfect/pinned-original-source.tex, source lines 5625--5625. Prerequisite pointer invoked at proof line(s) 6390; it names a standard fact of the same chapter that the counted argument uses, and every step of the counted proof built on it is retained verbatim in the delivered body.\par\medskip
\ifdefined\hypertarget\hypertarget{sourcepointer-2}{}\fi
\noindent\textit{Source reference.} Original reference 0A1H: sources/stacks-spaces-perfect/pinned-original-source.tex, source lines 6308--6308. Prerequisite pointer invoked at proof line(s) 6392; it names a standard fact of the same chapter that the counted argument uses, and every step of the counted proof built on it is retained verbatim in the delivered body.\par\medskip
\ifdefined\hypertarget\hypertarget{sourcepointer-3}{}\fi
\noindent\textit{Source reference.} Original reference 08D5: sources/stacks-spaces-perfect/pinned-original-source.tex, source lines 790--790. Prerequisite pointer invoked at proof line(s) 6402; it names a standard fact of the same chapter that the counted argument uses, and every step of the counted proof built on it is retained verbatim in the delivered body.\par\medskip
\ifdefined\hypertarget\hypertarget{sourcepointer-4}{}\fi
\noindent\textit{Source reference.} Original reference 08ES: sources/stacks-spaces-perfect/pinned-original-source.tex, source lines 3786--3786. Prerequisite pointer invoked at proof line(s) 6407; it names a standard fact of the same chapter that the counted argument uses, and every step of the counted proof built on it is retained verbatim in the delivered body.\par\medskip
\begin{lemma}
\label{lemma-base-change-tensor-pseudo-coherent}
Let $S$ be a scheme.
Let $f : X \to Y$ be a morphism of finite presentation
between algebraic spaces over $S$.
Let $E \in D(\mathcal{O}_X)$ be a pseudo-coherent object.
Let $\mathcal{G}^\bullet$ be a bounded above complex of
finitely presented $\mathcal{O}_X$-modules,
flat over $Y$, with support proper over $Y$. Then
$$
K = Rf_*(E \otimes_{\mathcal{O}_X}^\mathbf{L} \mathcal{G}^\bullet)
$$
is a pseudo-coherent object of $D(\mathcal{O}_Y)$ and its formation
commutes with arbitrary base change.
\end{lemma}

\begin{proof}
The statement on base change is Lemma \ref{lemma-base-change-tensor}.
Thus it suffices to show that $K$ is a pseudo-coherent object.
This will follow from Lemma \ref{lemma-base-change-tensor-perfect}
by approximation by perfect complexes. We encourage the reader to
skip the rest of the proof.

\medskip\noindent
The question is \'etale local on $Y$, hence we may assume $Y$ is affine.
Then $X$ is quasi-compact and quasi-separated. Moreover, there
exists an integer $N$ such that total direct image
$Rf_* : D_\QCoh(\mathcal{O}_X) \to D_\QCoh(\mathcal{O}_Y)$
has cohomological dimension $N$ as explained in
Lemma \ref{lemma-quasi-coherence-direct-image}.
Choose an integer $b$ such that $\mathcal{G}^i = 0$ for $i > b$.
It suffices to show that $K$ is $m$-pseudo-coherent for
every $m$. Choose an approximation $P \to E$ by a perfect complex $P$
of $(X, E, m - N - 1 - b)$. This is possible by
Theorem \ref{theorem-approximation}.
Choose a distinguished triangle
$$
P \to E \to C \to P[1]
$$
in $D_\QCoh(\mathcal{O}_X)$. The cohomology sheaves of $C$ are zero
in degrees $\geq m - N - 1 - b$. Hence
the cohomology sheaves of $C \otimes^\mathbf{L} \mathcal{G}^\bullet$
are zero in degrees $\geq m - N - 1$.
Thus the cohomology sheaves of $Rf_*(C \otimes^\mathbf{L} \mathcal{G})$
are zero in degrees $\geq m - 1$. Hence
$$
Rf_*(P \otimes^\mathbf{L} \mathcal{G}) \to
Rf_*(E \otimes^\mathbf{L} \mathcal{G})
$$
is an isomorphism on cohomology sheaves in degrees $\geq m$.
Next, suppose that $H^i(P) = 0$ for $i > a$. Then
$
P \otimes^\mathbf{L} \sigma_{\geq m - N - 1 - a}\mathcal{G}^\bullet
\longrightarrow
P \otimes^\mathbf{L} \mathcal{G}^\bullet
$
is an isomorphism on cohomology sheaves in degrees $\geq m - N - 1$.
Thus again we find that
$$
Rf_*(P \otimes^\mathbf{L} \sigma_{\geq m - N - 1 - a}\mathcal{G}^\bullet) \to
Rf_*(P \otimes^\mathbf{L} \mathcal{G}^\bullet)
$$
is an isomorphism on cohomology sheaves in degrees $\geq m$.
By Lemma \ref{lemma-base-change-tensor-perfect} the source
is a perfect complex.
We conclude that $K$ is $m$-pseudo-coherent as desired.
\end{proof}

\end{document}
